\documentclass[12pt,a4paper]{article}
\usepackage[utf8]{inputenc}
\usepackage[spanish]{babel}
\usepackage{geometry}
\usepackage{listings}
\usepackage{xcolor}
\usepackage{parskip}
\usepackage{amsmath}
\usepackage{hyperref}

\geometry{margin=2.5cm}

\lstset{
  basicstyle=\ttfamily\small,
  keywordstyle=\color{blue}\bfseries,
  commentstyle=\color{gray},
  stringstyle=\color{orange},
  breaklines=true,
  frame=single,
  language=Java,
  numbers=left,
  numberstyle=\tiny\color{gray},
  literate={á}{{\'a}}1 {é}{{\'e}}1 {í}{{\'\i}}1 {ó}{{\'o}}1 {ú}{{\'u}}1 {ñ}{{\~n}}1 {Á}{{\'A}}1 {É}{{\'E}}1 {Ó}{{\'O}}1
}

\title{\textbf{Gato (Tic-Tac-Toe) con Minimax}\\
  \large Documentación Técnica}
\author{Inteligencia Artificial}
\date{}

\begin{document}
\maketitle
\tableofcontents
\newpage

\section{Descripción del Proyecto}
Este proyecto implementa el juego clásico de \textbf{Gato (Tres en Raya)} con una inteligencia artificial basada en el algoritmo \textbf{Minimax}. La IA es teóricamente imbatible: siempre gana o empata contra cualquier jugador humano.

El juego se desarrolló en \textbf{Java} con una interfaz gráfica construida con \textbf{Swing}. El tablero es de 3×3 celdas y el jugador humano compite directamente contra la CPU.

\section{Lógica del Juego}

\subsection{Estado del tablero}
El tablero se representa como un arreglo de 9 enteros donde:
\begin{itemize}
  \item \texttt{0} = celda vacía
  \item \texttt{1} = jugador humano (X)
  \item \texttt{-1} = CPU (O)
\end{itemize}

\subsection{Condición de victoria}
El juego verifica las 8 combinaciones ganadoras posibles: 3 filas, 3 columnas y 2 diagonales. Si alguna de estas líneas tiene los 3 valores iguales (y distintos de 0), hay un ganador.

\section{Algoritmo Minimax}

\subsection{Concepto}
Minimax es un algoritmo de búsqueda adversarial para juegos de suma cero. Construye un \textbf{árbol de juego completo} desde el estado actual hasta todos los estados terminales, y asigna una puntuación a cada nodo:

\begin{itemize}
  \item \textbf{+10}: la CPU gana
  \item \textbf{-10}: el humano gana
  \item \textbf{0}: empate
\end{itemize}

La CPU (jugador MAX) elige el movimiento de mayor puntuación; el humano (jugador MIN) elegiría el de menor puntuación.

\subsection{Pseudocódigo}
\begin{lstlisting}
función minimax(tablero, profundidad, esMaximizador):
    si hayGanador(tablero, CPU):     retornar +10
    si hayGanador(tablero, humano):  retornar -10
    si tableroLleno(tablero):        retornar 0

    si esMaximizador:
        mejorValor = -infinito
        para cada celda vacía:
            tablero[celda] = CPU
            valor = minimax(tablero, prof+1, false)
            tablero[celda] = vacío
            mejorValor = max(mejorValor, valor)
        retornar mejorValor
    sino:
        mejorValor = +infinito
        para cada celda vacía:
            tablero[celda] = humano
            valor = minimax(tablero, prof+1, true)
            tablero[celda] = vacío
            mejorValor = min(mejorValor, valor)
        retornar mejorValor
\end{lstlisting}

\subsection{Elección del mejor movimiento}
Para determinar el movimiento óptimo de la CPU:

\begin{lstlisting}
función mejorMovimiento(tablero):
    mejorPuntuacion = -infinito
    mejorCelda = -1
    para cada celda vacía:
        tablero[celda] = CPU
        puntuacion = minimax(tablero, 0, false)
        tablero[celda] = vacío
        si puntuacion > mejorPuntuacion:
            mejorPuntuacion = puntuacion
            mejorCelda = celda
    retornar mejorCelda
\end{lstlisting}

\subsection{Análisis de complejidad}
El Gato tiene un espacio de estados máximo de $9! = 362{,}880$ tableros, pero con poda implícita se reduce significativamente. El árbol de juego completo tiene profundidad máxima 9 (9 celdas). La complejidad temporal es $O(b^d)$ donde $b$ es el factor de ramificación y $d$ la profundidad.

\section{Arquitectura del Código}

El proyecto se implementó en un único archivo \texttt{Gato\_SG\_GUI.java} que contiene:

\begin{itemize}
  \item \textbf{Clase principal}: extiende \texttt{JFrame} y contiene la GUI completa.
  \item \textbf{Método \texttt{minimax()}}: implementación recursiva del algoritmo.
  \item \textbf{Método \texttt{mejorMovimiento()}}: evalúa todas las celdas vacías y elige la óptima.
  \item \textbf{Método \texttt{verificarGanador()}}: comprueba las 8 combinaciones ganadoras.
  \item \textbf{ActionListeners}: detectan clic del jugador y disparan la respuesta de la CPU.
\end{itemize}

\section{Flujo de Ejecución}

\begin{enumerate}
  \item El jugador hace clic en una celda vacía del tablero.
  \item Se registra el movimiento del humano en el arreglo de estado.
  \item Se verifica si el humano ganó o si el tablero está lleno (empate).
  \item Si el juego continúa, se llama a \texttt{mejorMovimiento()} que ejecuta Minimax.
  \item La CPU coloca su ficha en la celda óptima.
  \item Se verifica si la CPU ganó o si hay empate.
  \item La interfaz se actualiza visualmente en cada paso.
\end{enumerate}

\section{Interfaz Gráfica}

La GUI fue construida con \textbf{Java Swing}:
\begin{itemize}
  \item \texttt{JFrame}: ventana principal del juego.
  \item \texttt{JButton[9]}: arreglo de 9 botones que representan las celdas del tablero.
  \item \texttt{JLabel}: muestra el estado del juego (turno, ganador, empate).
  \item \texttt{GridLayout(3,3)}: organiza los 9 botones en cuadrícula.
\end{itemize}

\section{Diferencia con Negamax}

Minimax usa dos funciones separadas (MAX y MIN). Negamax unifica ambas aprovechando que el valor de un estado para el jugador actual es el negativo del valor para el oponente:

\[f_{\text{negamax}}(s) = -f_{\text{negamax}}(\text{sucesor}(s))\]

Ambos algoritmos producen el mismo resultado; Negamax es más compacto en código.

\end{document}
